\section{Related Work}\label{sec:related}

\paragraph{GPU profiling and performance models.}
Nsight Systems provides individual CUDA events, kernel summaries, and
reports that group kernels by NVTX range
\cite{nvidiaNsightSystems}. PyTorch Profiler correlates GPU activities
with framework events, and Holistic Trace Analysis uses those traces
to analyze distributed workloads \cite{pytorchProfiler,pytorchHta}.
Nsight Compute provides hardware counters and kernel rooflines. Its
replay mode and cache controls affect the interpretation of measurements
that depend on preceding launches \cite{nvidiaNsightCompute}. Operator
Profiler associates these hardware counters with PyTorch operators
\cite{chu2026operator}. PyTorch also supports semantic labels on graph
nodes that are visible in replay traces \cite{pytorchGraphAnnotations}.
\sys{} joins timing records with kernel arguments and a model memory
map, making it possible to compare calls to the same kernel by the
tensor each call references.

The Roofline model bounds performance using operational intensity,
compute throughput, and memory bandwidth \cite{williams2009roofline}.
BabelStream measures attainable bandwidth on streaming workloads
\cite{deakin2016babelstream}. AMD TraceLens reconstructs hierarchical
timing breakdowns from profiler traces and uses operator models to
estimate FLOPs and memory traffic. For attention rooflines, it can use
per-step workload statistics from vLLM's optional trace annotations
\cite{amdTracelens,vllmRooflineAnnotations}. StriaTrace fits empirical
latency boundaries to detect anomalous serving steps and correlates
trace data to help diagnose them \cite{wu2026striatrace}. AgentPerfBench
collects per-kernel Nsight Compute profiles of agentic workloads and
applies multidimensional roofline analysis to representative kernels
\cite{lau2026agentperfbench}.
\sys{} uses a measured read-bandwidth baseline and a weight-size
estimate to guide controlled comparisons between launches with
different tensor identities or preceding work.

\paragraph{Tracing and memory attribution.}
Mystique generates representative benchmarks by replaying operator-level
PyTorch execution traces. Chakra standardizes operation graphs for
benchmarking and system co-design
\cite{liang2023mystique,sridharan2023chakra}. To reduce LLM serving
cold starts, Medusa prepares CUDA graphs and KV-cache initialization
metadata offline, and Foundry saves graph topology and execution
context for template-based reconstruction
\cite{zeng2025medusa,liu2026foundry}. \sys{} uses CUDA
graph node metadata to interpret timings in a steady-state trace.
NVBit supports instruction-level memory reference tracing
\cite{villa2019nvbit}. Von Buelow et al. use instrumented memory
references to simulate L1 and L2 behavior and estimate hit rates per
allocation \cite{buelow2022finegrained}. DrGPUM correlates memory
inefficiencies such as redundant or unused allocations with data objects
and GPU API calls
\cite{lin2023drgpum}. \sys{} uses launch arguments as a proxy for
tensor references and measures the timing of the launches in the
running engine.

\paragraph{GPU interference across kernels.}
Elvinger et al. characterize L2 cache pollution and bandwidth
contention between colocated GPU workloads. Mosaic predicts the impact
of kernel colocation for latency-aware scheduling, and CacheC manages
L2 to improve kernel concurrency
\cite{elvinger2025gpu,strati2026mosaic,xi2026cachec}.
In GTX 1080 experiments, Yandrofski et al. observe larger victim-kernel
slowdowns with concurrent L2 write workloads than with read workloads.
They discuss dirty-line write-back as one possible contributor
\cite{yandrofski2022enemies}. NVIDIA documents that an L2 miss can
generate a DRAM write when it evicts a dirty line
\cite{nvidiaComputeTriage}.
Our case study examines the timing cost of this mechanism between
successive kernels in a decode step.
Matching the projection launches exposes the timing difference, and
varying the preceding writes in isolation tests the explanation.

\paragraph{Speculative decoding with recurrent state.}
Speculative decoding verifies several candidate tokens at once, and
multi-token prediction can provide those candidates
\cite{leviathan2023fast,deepseekv3}. Gated DeltaNet maintains a recurrent
state \cite{yang2025gated}. Its speculative implementation must recover
the state at the accepted position. Prior vLLM changes proposed deferring
intermediate state stores in the Triton implementation
\cite{vllmpr61434,vllmpr64325}. Our case study traces
the excess stores to both their originating kernel and the following
projection, then evaluates a CUDA implementation that shares the
Triton record layout \cite{vllmpr69160}.
